\documentclass[11pt]{article}
\usepackage[utf8]{inputenc}\usepackage[T1]{fontenc}\usepackage{mathptmx}
\usepackage[english]{babel}
\usepackage[a4paper,margin=2.2cm]{geometry}
\setlength{\parindent}{0pt}\setlength{\parskip}{7pt}\linespread{1.05}
\pagestyle{empty}
\begin{document}
\hfill Lima, 13 de noviembre de 2026

\begin{center}\large\bfseries Carta de presentación\end{center}

Estimado editor:

Nos complace presentar nuestro manuscrito titulado “Machine Learning Model for Predicting the Risk of Childhood Stunting in Peru” (Modelo de machine learning para predecir el riesgo de desnutrición crónica infantil en el Perú) para su consideración como artículo de investigación original en \emph{Nutrition}. Los autores son Abel Is-Boset Guevara-Huasco, Zulin Pamela Vallejos-Cotrina, Verónica Lisseth Vergara-Rojas, Ferdinand Edgardo Pineda-Ancco y Nemias Saboya-Rios.

El estudio desarrolla y valida internamente un modelo para predecir la desnutrición crónica en 304,271 niños menores de cinco años, a partir de la Encuesta Demográfica y de Salud Familiar (ENDES) 2007–2025 y de 45 predictores de cinco dominios: el niño, la madre, el hogar y el territorio, el agua y el saneamiento, y la vivienda. Se compararon seis algoritmos supervisados mediante validación cruzada estratificada de cinco pliegues con ponderación de la encuesta. Se seleccionó LightGBM, con un AUC-ROC fuera de pliegue de 0.829 (IC 95\%: 0.827–0.832) y una sensibilidad del 76.4\%. Los valores SHAP identificaron la pobreza del hogar, el peso al nacer y la altitud como los principales predictores y mostraron cómo su peso varía entre los 25 departamentos del Perú.

El modelo se implementó en una aplicación web de apoyo a la decisión, que un panel de cuatro profesionales de salud valoró favorablemente (mediana de la valoración global: 9 sobre 10). El código, el conjunto de datos procesado y el modelo entrenado están disponibles en Zenodo para favorecer la reproducibilidad.

Consideramos que el manuscrito interesará a los lectores de \emph{Nutrition}, que ha publicado varios estudios de machine learning sobre desnutrición infantil en Bangladesh, Egipto, India, Filipinas y África subsahariana. Nuestro trabajo extiende esta línea de investigación a América Latina, donde aún es escasa, y aporta una interpretación territorial y una herramienta práctica dirigida a las instituciones de salud pública.

Confirmamos que este manuscrito no ha sido publicado en otro lugar y no está en evaluación en ninguna otra revista. Todos los autores han aprobado el manuscrito y están de acuerdo con su envío a \emph{Nutrition}. El estudio fue aprobado por el Comité de Ética de la Escuela Profesional de Ingeniería de Sistemas de la Universidad Peruana Unión (N.° CEI-UPEU-FIA-2026-24), y los autores declaran no tener conflictos de intereses.

Agradecemos su tiempo y consideración.

Atentamente,

\vspace{1.2cm}
{\setlength{\parskip}{0pt}\small \textbf{Mg. Nemias Saboya-Rios}\\Autor de correspondencia\\Escuela Profesional de Ingeniería de Sistemas, Facultad de Ingeniería y Arquitectura\\Universidad Peruana Unión, Carretera Central Km 19.5, Ñaña, Lurigancho-Chosica, Lima, Perú\\saboya@upeu.edu.pe · ORCID 0000-0002-7166-2197\par}
\end{document}
